% ============================================================
% Appendix E — Institutional Context
% ============================================================
% USAGE: % ============================================================
% Appendix E — Institutional Context
% ============================================================
% USAGE: % ============================================================
% Appendix E — Institutional Context
% ============================================================
% USAGE: % ============================================================
% Appendix E — Institutional Context
% ============================================================
% USAGE: \input{Appendix_E} from your main .tex file.
%
% REQUIRED in your main preamble (most likely already loaded):
%   \usepackage{booktabs}
%   \usepackage{array}
%   \usepackage{hyperref}
% ============================================================

\renewcommand{\thetable}{E.\arabic{table}}
\setcounter{table}{13}
\setcounter{subsection}{0}

\section{Institutional context and comparison: quantitative summary}
\addcontentsline{toc}{section}{Appendix E: Institutional Context}

This appendix documents the factual institutional context underlying the interpretive discussion in Section~5. It provides: (i)~a side-by-side chronology of the formal institutional responses, (ii)~a summary table of quantitative institutional contrasts, (iii)~a regulatory milestone timeline, and (iv)~an interpretive integration of the Table~\ref{tab:fiscal_inst} with the disaster-specific evidence. Importantly, these are descriptive inputs, not causal estimates; they serve to ground the mechanism discussion in observable facts.


% ============================================================
% E.1   COMPARATIVE INSTITUTIONAL CHRONOLOGY
% ============================================================
\subsection{Comparative Institutional Chronology}

Table~\ref{tab:chronology} presents a side-by-side timeline of key institutional actions.

%%% TABLE E.14 — Part (a) %%%
\begin{table}[!htbp]
\centering
\caption{Side-by-side institutional response chronology}
\label{tab:chronology}
\footnotesize
\begin{tabular}{@{}
  >{\raggedright\arraybackslash}p{0.12\textwidth}
  >{\raggedright\arraybackslash}p{0.40\textwidth}
  >{\raggedright\arraybackslash}p{0.40\textwidth}
@{}}
\toprule
\textbf{Timeframe} &
\textbf{Chile (Maule, Feb 2010)} &
\textbf{New Zealand (Canterbury, Sep 2010--Feb 2011)} \\
\midrule

Days 1--7} &
State of Catastrophe declared in Maule and Biob{\'{\i}}o (Feb~28). Military curfew imposed. Approx.\ 14,000 troops deployed \parencite{CRS2010}. ONEMI coordinated initial response; later criticized for tsunami warning failures. &
Sep 2010 (Darfield): State of Emergency declared. Civil Defence managed initial response. No significant civil disorder. EQC activated Catastrophe Response Plan; claims lodging opened within days. \\
\addlinespace

{Weeks 2--8} &
Presidential transition (Bachelet to Pi{\~n}era, Mar~11) complicated coordination. National reconstruction coordinator appointed at MINVU. Initial damage surveys began. No centralized claims system for the approx.\ 10--20\% of residentially insured properties. &
EQC appointed Fletcher Construction to project-manage residential repairs (approx.\ 50,000 homes). Assessment teams deployed. Small claims (under-cap) assessed within months \parencite{Khakurel2021}. Private insurers processing in parallel. \\
\addlinespace

{Months 2--6} &
Plan de Reconstrucci{\'o}n published by MINVU (Aug~2010), six months post-quake. Targeted 220,000 housing units over four years. Finance plan adopted. No dedicated recovery authority created. &
Feb 2011 Christchurch earthquake (Mw~6.3, 185 deaths). CERA established 29~Mar 2011---35 days after the Feb event. Wide-ranging powers including authority to suspend laws. Canterbury Earthquake Recovery Act passed (Apr~2011). \\

\bottomrule
\end{tabular}

\vspace{2pt}
{\scriptsize\textit{Continued in next panel.}}
\end{table}

%%% TABLE E.14 — Part (b) %%%
\begin{table}[!htbp]
\centering
\footnotesize
\begin{tabular}{@{}
  >{\raggedright\arraybackslash}p{0.12\textwidth}
  >{\raggedright\arraybackslash}p{0.40\textwidth}
  >{\raggedright\arraybackslash}p{0.40\textwidth}
@{}}
\toprule
\textbf{Timeframe} &
\textbf{Chile (Maule, Feb 2010)} &
\textbf{New Zealand (Canterbury, Sep 2010--Feb 2011)} \\
\midrule

{Months 6--12} &
Reconstruction managed through existing ministries (MINVU, MOP). No new entity created. Some temporary shelters operational. &
Red Zone announced; purchase offers to approx.\ 8,060 properties. Technical categories (TC1--TC3) established. Building consent fast-tracking initiated. Fletcher EQR expanded to approx.\ 70,000 homes. \\
\addlinespace

{Year 2} &
Centralized ministerial direction continued. No significant regulatory changes. Community participation limited \parencite{CardenasJiron2012}. By Oct~2012: 54\% housing complete; 30\% under construction \parencite{Platt2019}. &
Canterbury Recovery Strategy approved (May~2012). SCIRT alliance formed (CERA + CCC + NZTA + 5 firms). Land Use Recovery Plan approved (2013). Community groups active. Peak construction began. \\
\addlinespace

{Years 3--5} &
Spending continued through central transfers. No comparable construction sector surge relative to counterfactual (Section~\ref{sec:sectoral_impact_analysis}). &
CERA spent NZ\$4bn over 5~years. Construction share peaked ${\sim}$10\% in 2016. CERA disestablished 18~Apr 2016; functions transferred to DPMC, LINZ, MBIE, CCC, Regenerate Christchurch, {\=O}takaro Ltd. \\

\bottomrule
\end{tabular}

\vspace{2pt}
{\scriptsize \textit{Table~\ref{tab:chronology} concluded. Sources:} Chile: \parencite{eclac2010overview,Platt2019,Comerio2013,CRS2010}. New Zealand: \parencite{AuditorGeneral2017,CERAA2011,ICNZ2024,eqc2017annual,grube2018embedded,Wilson2016}.}
\end{table}


% ============================================================
% E.2   QUANTITATIVE INSTITUTIONAL CONTRASTS
% ============================================================
\subsection{Quantitative Institutional Contrasts}

Table~\ref{tab:quantcontrasts} summarizes the key quantitative differences in institutional architecture available for post-disaster recovery.

%%% TABLE E.15 — Part (a): Insurance + Governance %%%
\begin{table}[!htbp]
\centering
\caption{Key quantitative institutional contrasts: Chile vs.\ New Zealand}
\label{tab:quantcontrasts}
\footnotesize
\begin{tabular}{@{}
  >{\raggedright\arraybackslash}p{0.29\textwidth}
  >{\raggedright\arraybackslash}p{0.30\textwidth}
  >{\raggedright\arraybackslash}p{0.33\textwidth}
@{}}
\toprule
\textbf{Variable} &
\textbf{Chile (Maule)} &
\textbf{New Zealand (Canterbury)} \\
\midrule

\multicolumn{3}{@{}l}{\textit{\textbf{A. Insurance architecture}}} \\
\addlinespace[2pt]

Residential EQ insurance penetration &
${\sim}$10--20\% &
Near-universal (EQC automatic with fire ins.) \\

Commercial insurance penetration &
${\sim}$60\% &
High (bundled) \\

Pre-existing public insurance scheme &
None &
EQC (est.\ 1945) \\

Pre-quake accumulated fund &
n/a &
NZ\$6.1 billion \parencite{eqc2017annual} \\

Total insured cost &
US\$4--7 billion &
NZ\$31 billion \parencite{ICNZ2024} \\

EQC claims filed &
n/a &
$>$470,000 \\

Intl.\ reinsurance recovered &
n/a &
$>$NZ\$5 billion \parencite{NHC2024} \\

Median settlement time (EQC) &
No comparable public data &
Small: 138--191\,d; Med: 712--890\,d; Large: 1,226--1,392\,d \parencite{Khakurel2021} \\

Crown guarantee for shortfall &
None &
Yes (EQC Act s.\ 16) \\

\addlinespace[4pt]
\multicolumn{3}{@{}l}{\textit{\textbf{B. Governance and recovery institutions}}} \\
\addlinespace[2pt]

Dedicated recovery authority &
None created &
CERA (est.\ 29 Mar 2011; disest.\ 18 Apr 2016) \\

Time to recovery authority &
--- &
35 days \\

Special legislative powers &
None enacted &
Canterbury EQ Recovery Act 2011 \\

Regulatory fast-tracking &
No significant changes &
Yes (technical categories, consent fast-tracking) \\

Public--private rebuild alliance &
No &
SCIRT (CCC + CERA + NZTA + 5 firms) \\

Recovery authority spending (5~yrs) &
--- &
NZ\$4 billion \parencite{AuditorGeneral2017} \\


\end{tabular}


%%% TABLE E.15 — Part (b): Spending + Social capital %%%

\centering
\footnotesize
\begin{tabular}{@{}
  >{\raggedright\arraybackslash}p{0.29\textwidth}
  >{\raggedright\arraybackslash}p{0.30\textwidth}
  >{\raggedright\arraybackslash}p{0.33\textwidth}
@{}}
\toprule
\textbf{Variable} &
\textbf{Chile (Maule)} &
\textbf{New Zealand (Canterbury)} \\
\midrule

\multicolumn{3}{@{}l}{\textit{\textbf{C. Public reconstruction spending}}} \\
\addlinespace[2pt]

Govt reconstruction commitment &
US\$8.4 billion &
NZ\$10.1 billion (incl.\ CERF + Natural Hazard Fund) \\

Timing of finance plan &
${\sim}$6 months (Aug 2010) &
${\sim}$3 months (May 2011 Budget) \\

Housing program scale &
220,000 units targeted &
${\sim}$70,000 managed repairs + Red Zone buyouts \\

\addlinespace[4pt]
\multicolumn{3}{@{}l}{\textit{\textbf{D. Informal institutions and social capital}}} \\
\addlinespace[2pt]

Pre-quake social trust &
Low \parencite{dussaillant2014trust} &
Moderate-to-high \\

Post-disaster civil disorder &
Moderate to high (post-disaster civil disorder and military curfew) &
Minimal \\

Community co-production &
Limited formal engagement &
Substantial (volunteers, entrepreneurs) \parencite{grube2018embedded} \\

\bottomrule
\end{tabular}

\vspace{4pt}
{\scriptsize \textit{Table~\ref{tab:quantcontrasts} concluded.} \textbf{Notes:} Insurance penetration (Chile): \parencite{AIRWorldwide2010}; \textcite{NguyenNoy2020} report residential coverage below 20\%. New~Zealand: near-universal, with roughly 80\% insured \parencite{NguyenNoy2020}. Reinsurance (NZ): \parencite{NHC2024}. CERA spending: \parencite{AuditorGeneral2017}. Red Zone buyouts: \parencite{Wilson2016}. Social capital: \parencite{dussaillant2014trust}. Settlement times: \parencite{Khakurel2021}. Over-cap transfers: \parencite{ICNZ2024}.}
\end{table}


% ============================================================
% E.3   REGULATORY AND LEGISLATIVE MILESTONE TIMELINE
% ============================================================
\subsection{Regulatory and Legislative Milestone Timeline}

Table~\ref{tab:regmilestones} provides a month-by-month comparison of the regulatory and legislative milestones in each country's disaster recovery. The contrast in institutional readiness is pronounced: New Zealand enacted emergency recovery legislation within 10~days of the September 2010 event, established CERA within 35~days of the February 2011 shock, and had gazetted technical land categories, building consent fast-tracking, and a public--private infrastructure alliance within 18~months. In contrast, Chile's response operated through pre-existing ministerial structures \textit{without} creating a dedicated recovery entity, enacting special legislative powers, or establishing formal regulatory fast-tracking for reconstruction.

%%% TABLE E.16 — Part (a): Month 0 through 8--12 %%%
\begin{table}[!htbp]
\centering
\caption{Regulatory and legislative milestone timeline (months from first event)}
\label{tab:regmilestones}
\footnotesize
\begin{tabular}{@{}
  >{\raggedright\arraybackslash}p{0.06\textwidth}
  >{\raggedright\arraybackslash}p{0.12\textwidth}
  >{\raggedright\arraybackslash}p{0.34\textwidth}
  >{\raggedright\arraybackslash}p{0.38\textwidth}
@{}}
\toprule
\textbf{Mo.} &
\textbf{Date} &
\textbf{Chile} &
\textbf{New Zealand} \\
\midrule

{0} & Feb--Sep 2010 &
State of Catastrophe (Feb~28). ONEMI response. No new legislation. &
Darfield earthquake (Sep~4). State of Emergency. Canterbury EQ Response and Recovery Act 2010 enacted (Sep~14). \\
\addlinespace

{1--2} & Oct--Nov 2010 &
Damage assessment ongoing. No new legislation. &
Canterbury EQ Recovery Commission established. EQC Catastrophe Response Plan operational. Fletcher Construction appointed. \\
\addlinespace

{3--6} & Dec 2010--Feb 2011 &
Plan published Aug 2010; no further legislative action. &
Feb~22 Christchurch earthquake (Mw~6.3, 185 deaths). National state of emergency declared. \\
\addlinespace

{6--7} & Mar--Apr 2011 &
Reconstruction through MINVU and existing ministers. No dedicated authority. &
CERA established 29 Mar 2011 (Order in Council). Canterbury EQ Recovery Act 2011 enacted (Apr). Powers to suspend/amend regulations. \\
\addlinespace

{8--12} & May--Sep 2011 &
Housing subsidies via SERVIUs. No regulatory fast-tracking. &
Red Zone announced (Jun). Technical categories (TC1--TC3) gazetted. Consent fast-tracking. 8,060 properties offered buyouts. \\

\end{tabular}

\centering
\footnotesize
\begin{tabular}{@{}
  >{\raggedright\arraybackslash}p{0.06\textwidth}
  >{\raggedright\arraybackslash}p{0.12\textwidth}
  >{\raggedright\arraybackslash}p{0.34\textwidth}
  >{\raggedright\arraybackslash}p{0.38\textwidth}
@{}}
\toprule
\textbf{Mo.} &
\textbf{Date} &
\textbf{Chile} &
\textbf{New Zealand} \\
\midrule

{12--18} & Oct 2011--Mar 2012 &
54\% of housing underway or complete by Oct 2012. No new institutional changes. &
EQC disputes-resolution manager appointed (Mar~2012). CCDU established within CERA (Apr~2012) for CBD recovery plan. \\
\addlinespace

{18--24} & Apr--Sep 2012 &
Ongoing centralized delivery. No public--private alliance formed. &
Canterbury Recovery Strategy approved (May~2012). SCIRT alliance operational. Budget ${\sim}$NZ\$2.5bn (actual: NZ\$2.2bn). \\
\addlinespace

{24--36} & Oct 2012--Sep 2013 &
54\% housing complete; 30\% under construction (Oct 2012). &
Land Use Recovery Plan approved (2013). Resilience integrated into district plans. Over 60 Orders in Council amending legislation. \\
\addlinespace

{36--60} & 2013--2016 &
Continued central govt transfers. No comparable construction surge. &
Peak construction 2015--2016. CERA spent NZ\$4bn. CERA disestablished 18~Apr 2016; functions transferred. \\

\bottomrule
\end{tabular}

\vspace{2pt}
{\scriptsize \textit{Table~\ref{tab:regmilestones} concluded. Sources:} Chile: \parencite{Comerio2013,Platt2019}. New Zealand: \parencite{CERRA2010,CERAA2011,AuditorGeneral2017,Wilson2016}.}
\end{table}


% ============================================================
% E.4   INTEGRATION OF MACRO-INSTITUTIONAL INDICATORS
% ============================================================
\subsection{Integration of Macro-Institutional Indicators (Table~\ref{tab:fiscal_inst}) with Disaster-Specific Evidence}

Table~\ref{tab:fiscal_inst} in the main text reports pre-quake national-level fiscal and institutional indicators for both countries. As noted in Section~5, these national-level scores serve as contextual descriptors rather than causal inputs; sub-national institutional quality indices are not available for the affected regions. Nevertheless, interpreting Table~\ref{tab:fiscal_inst} alongside the disaster-specific institutional evidence documented above yields several interpretively useful contrasts.

\textbf{Fiscal capacity and the insurance--debt inversion:} The most striking contrast in Table~\ref{tab:fiscal_inst} is between government debt (Chile 11.1\% of GDP; New Zealand 62.6\%) and the insurance architecture documented in Table~\ref{tab:quantcontrasts}. Chile entered the disaster with greater fiscal space, yet Maule's GDP trajectory shows no positive overshoot. New Zealand, despite higher sovereign indebtedness, channeled reconstruction capital primarily through its pre-existing insurance system rather than fiscal expansion. The magnitudes involved are telling: NZ's total insured cost reached NZ\$31~billion, of which more than NZ\$5~billion flowed from international reinsurers (Table~\ref{tab:quantcontrasts}), an exogenous capital injection that required no fiscal outlay. Chile's total insured cost was only US\$4--7~billion, with residential penetration below 20\%, leaving a US\$8.4~billion fiscal commitment as the primary reconstruction channel. Accordingly, Chile's change in public spending in the quake year (+10.44\%) exceeded New Zealand's (+6.24\%), directly reflecting this heavier reliance on the fiscal channel. The chronology in Table~\ref{tab:chronology} reinforces the point: NZ's EQC activated its claims process within days and appointed Fletcher Construction within weeks, whereas Chile's Plan de Reconstrucci{\'o}n was published six months after the earthquake. This contrast suggests that what matters for recovery dynamics is not necessarily raw fiscal capacity \textit{per se}, but the institutional and financial architecture through which reconstruction capital is mobilized. This observation does not imply that fiscal capacity is irrelevant, but rather that the \textit{speed and institutional channel} of financing may matter more for regional output than aggregate fiscal capacity.

\textbf{Legal system, property rights, and regulatory quality:} New Zealand scores higher on Legal System \& Property Rights (8.69 vs.\ 6.80) and Regulation (8.92 vs.\ 7.55). These national-level scores are consistent with the disaster-specific institutional evidence documented in Tables~\ref{tab:chronology} and~\ref{tab:regmilestones}: New Zealand enacted the Canterbury Earthquake Response and Recovery Act within 10~days of the September 2010 event, established CERA with statutory override powers within 35~days of the February 2011 shock, gazetted technical land categories and building consent fast-tracking within 12~months, and had the SCIRT public--private alliance and Land Use Recovery Plan operational within 24~months. Over CERA's lifespan, more than 60 Orders in Council were made to amend existing legislation for recovery purposes (Table~\ref{tab:regmilestones}). Chile, by contrast, made no comparable regulatory adaptations: reconstruction was managed through existing MINVU and MOP structures without a dedicated recovery entity, special legislative powers, or consent fast-tracking. The higher scores on Legal System \& Property Rights also align with the insurance channel: a well-functioning property rights regime is a \textit{precondition} for the complex contractual architecture of catastrophe insurance, reinsurance treaties, and the claims adjudication process through which more than 470,000 EQC and 168,000 private insurer claims were processed.

\textbf{Size of Government:} Chile's higher score (7.91 vs.\ 4.91) indicates a smaller government footprint relative to GDP (lower taxation, less government consumption, and fewer transfers). This is consistent with Chile's reliance on centralized but \textit{ad hoc} fiscal programs for reconstruction: the US\$8.4~billion commitment was channeled through existing ministry structures with no new institutional entities, whereas New Zealand's larger public sector provided the institutional infrastructure through which recovery was coordinated. Specifically, EQC (a Crown entity established in 1945 with NZ\$6.1~billion in accumulated reserves by 2010), CERA (a new government department with NZ\$4~billion in spending over five years), and the SCIRT alliance (a public--private arrangement coordinating three government entities and five construction firms on a NZ\$2.2~billion programme) were all creatures of New Zealand's comparatively larger state apparatus. The contrast illustrates that the relevant dimension is not government size in the abstract but rather the pre-existing institutional architecture that can be activated and adapted in a crisis. Chile's smaller government had correspondingly less pre-existing institutional capacity for disaster coordination (no public insurance scheme, no dedicated recovery authority) and the six-month lag to the Plan de Reconstrucci{\'o}n (Table~\ref{tab:chronology}) reflects the time required to build such capacity from scratch within existing ministerial structures.

To conclude, interpreting Table~\ref{tab:fiscal_inst} alongside the disaster-specific evidence in Tables~\ref{tab:chronology} to \ref{tab:regmilestones} reinforces the core interpretive point of Section~5: the divergent recovery trajectories are most consistently explained not by differential fiscal capacity (which actually favored Chile, at 11\% vs.\ 63\% debt-to-GDP) nor by broad macroeconomic stability, but by the pre-existing institutional architecture for disaster-specific financial mobilization, regulatory adaptation, and coordinated implementation. Chile's superior fiscal position did not translate into a comparable output response because the institutional channel through which spending was deployed (i.e., centralized ministerial programs without a dedicated recovery authority, or insurance-based external capital) was slower and less regionally targeted. The six-month lag to the Plan de Reconstrucci{\'o}n, the absence of regulatory override powers, and the limited community co-production all contrast with New Zealand's 35-day CERA establishment, over 60 Orders in Council adapting legislation, and the SCIRT public--private alliance.
 from your main .tex file.
%
% REQUIRED in your main preamble (most likely already loaded):
%   \usepackage{booktabs}
%   \usepackage{array}
%   \usepackage{hyperref}
% ============================================================

\renewcommand{\thetable}{E.\arabic{table}}
\setcounter{table}{13}
\setcounter{subsection}{0}

\section{Institutional context and comparison: quantitative summary}
\addcontentsline{toc}{section}{Appendix E: Institutional Context}

This appendix documents the factual institutional context underlying the interpretive discussion in Section~5. It provides: (i)~a side-by-side chronology of the formal institutional responses, (ii)~a summary table of quantitative institutional contrasts, (iii)~a regulatory milestone timeline, and (iv)~an interpretive integration of the Table~\ref{tab:fiscal_inst} with the disaster-specific evidence. Importantly, these are descriptive inputs, not causal estimates; they serve to ground the mechanism discussion in observable facts.


% ============================================================
% E.1   COMPARATIVE INSTITUTIONAL CHRONOLOGY
% ============================================================
\subsection{Comparative Institutional Chronology}

Table~\ref{tab:chronology} presents a side-by-side timeline of key institutional actions.

%%% TABLE E.14 — Part (a) %%%
\begin{table}[!htbp]
\centering
\caption{Side-by-side institutional response chronology}
\label{tab:chronology}
\footnotesize
\begin{tabular}{@{}
  >{\raggedright\arraybackslash}p{0.12\textwidth}
  >{\raggedright\arraybackslash}p{0.40\textwidth}
  >{\raggedright\arraybackslash}p{0.40\textwidth}
@{}}
\toprule
\textbf{Timeframe} &
\textbf{Chile (Maule, Feb 2010)} &
\textbf{New Zealand (Canterbury, Sep 2010--Feb 2011)} \\
\midrule

Days 1--7} &
State of Catastrophe declared in Maule and Biob{\'{\i}}o (Feb~28). Military curfew imposed. Approx.\ 14,000 troops deployed \parencite{CRS2010}. ONEMI coordinated initial response; later criticized for tsunami warning failures. &
Sep 2010 (Darfield): State of Emergency declared. Civil Defence managed initial response. No significant civil disorder. EQC activated Catastrophe Response Plan; claims lodging opened within days. \\
\addlinespace

{Weeks 2--8} &
Presidential transition (Bachelet to Pi{\~n}era, Mar~11) complicated coordination. National reconstruction coordinator appointed at MINVU. Initial damage surveys began. No centralized claims system for the approx.\ 10--20\% of residentially insured properties. &
EQC appointed Fletcher Construction to project-manage residential repairs (approx.\ 50,000 homes). Assessment teams deployed. Small claims (under-cap) assessed within months \parencite{Khakurel2021}. Private insurers processing in parallel. \\
\addlinespace

{Months 2--6} &
Plan de Reconstrucci{\'o}n published by MINVU (Aug~2010), six months post-quake. Targeted 220,000 housing units over four years. Finance plan adopted. No dedicated recovery authority created. &
Feb 2011 Christchurch earthquake (Mw~6.3, 185 deaths). CERA established 29~Mar 2011---35 days after the Feb event. Wide-ranging powers including authority to suspend laws. Canterbury Earthquake Recovery Act passed (Apr~2011). \\

\bottomrule
\end{tabular}

\vspace{2pt}
{\scriptsize\textit{Continued in next panel.}}
\end{table}

%%% TABLE E.14 — Part (b) %%%
\begin{table}[!htbp]
\centering
\footnotesize
\begin{tabular}{@{}
  >{\raggedright\arraybackslash}p{0.12\textwidth}
  >{\raggedright\arraybackslash}p{0.40\textwidth}
  >{\raggedright\arraybackslash}p{0.40\textwidth}
@{}}
\toprule
\textbf{Timeframe} &
\textbf{Chile (Maule, Feb 2010)} &
\textbf{New Zealand (Canterbury, Sep 2010--Feb 2011)} \\
\midrule

{Months 6--12} &
Reconstruction managed through existing ministries (MINVU, MOP). No new entity created. Some temporary shelters operational. &
Red Zone announced; purchase offers to approx.\ 8,060 properties. Technical categories (TC1--TC3) established. Building consent fast-tracking initiated. Fletcher EQR expanded to approx.\ 70,000 homes. \\
\addlinespace

{Year 2} &
Centralized ministerial direction continued. No significant regulatory changes. Community participation limited \parencite{CardenasJiron2012}. By Oct~2012: 54\% housing complete; 30\% under construction \parencite{Platt2019}. &
Canterbury Recovery Strategy approved (May~2012). SCIRT alliance formed (CERA + CCC + NZTA + 5 firms). Land Use Recovery Plan approved (2013). Community groups active. Peak construction began. \\
\addlinespace

{Years 3--5} &
Spending continued through central transfers. No comparable construction sector surge relative to counterfactual (Section~\ref{sec:sectoral_impact_analysis}). &
CERA spent NZ\$4bn over 5~years. Construction share peaked ${\sim}$10\% in 2016. CERA disestablished 18~Apr 2016; functions transferred to DPMC, LINZ, MBIE, CCC, Regenerate Christchurch, {\=O}takaro Ltd. \\

\bottomrule
\end{tabular}

\vspace{2pt}
{\scriptsize \textit{Table~\ref{tab:chronology} concluded. Sources:} Chile: \parencite{eclac2010overview,Platt2019,Comerio2013,CRS2010}. New Zealand: \parencite{AuditorGeneral2017,CERAA2011,ICNZ2024,eqc2017annual,grube2018embedded,Wilson2016}.}
\end{table}


% ============================================================
% E.2   QUANTITATIVE INSTITUTIONAL CONTRASTS
% ============================================================
\subsection{Quantitative Institutional Contrasts}

Table~\ref{tab:quantcontrasts} summarizes the key quantitative differences in institutional architecture available for post-disaster recovery.

%%% TABLE E.15 — Part (a): Insurance + Governance %%%
\begin{table}[!htbp]
\centering
\caption{Key quantitative institutional contrasts: Chile vs.\ New Zealand}
\label{tab:quantcontrasts}
\footnotesize
\begin{tabular}{@{}
  >{\raggedright\arraybackslash}p{0.29\textwidth}
  >{\raggedright\arraybackslash}p{0.30\textwidth}
  >{\raggedright\arraybackslash}p{0.33\textwidth}
@{}}
\toprule
\textbf{Variable} &
\textbf{Chile (Maule)} &
\textbf{New Zealand (Canterbury)} \\
\midrule

\multicolumn{3}{@{}l}{\textit{\textbf{A. Insurance architecture}}} \\
\addlinespace[2pt]

Residential EQ insurance penetration &
${\sim}$10--20\% &
Near-universal (EQC automatic with fire ins.) \\

Commercial insurance penetration &
${\sim}$60\% &
High (bundled) \\

Pre-existing public insurance scheme &
None &
EQC (est.\ 1945) \\

Pre-quake accumulated fund &
n/a &
NZ\$6.1 billion \parencite{eqc2017annual} \\

Total insured cost &
US\$4--7 billion &
NZ\$31 billion \parencite{ICNZ2024} \\

EQC claims filed &
n/a &
$>$470,000 \\

Intl.\ reinsurance recovered &
n/a &
$>$NZ\$5 billion \parencite{NHC2024} \\

Median settlement time (EQC) &
No comparable public data &
Small: 138--191\,d; Med: 712--890\,d; Large: 1,226--1,392\,d \parencite{Khakurel2021} \\

Crown guarantee for shortfall &
None &
Yes (EQC Act s.\ 16) \\

\addlinespace[4pt]
\multicolumn{3}{@{}l}{\textit{\textbf{B. Governance and recovery institutions}}} \\
\addlinespace[2pt]

Dedicated recovery authority &
None created &
CERA (est.\ 29 Mar 2011; disest.\ 18 Apr 2016) \\

Time to recovery authority &
--- &
35 days \\

Special legislative powers &
None enacted &
Canterbury EQ Recovery Act 2011 \\

Regulatory fast-tracking &
No significant changes &
Yes (technical categories, consent fast-tracking) \\

Public--private rebuild alliance &
No &
SCIRT (CCC + CERA + NZTA + 5 firms) \\

Recovery authority spending (5~yrs) &
--- &
NZ\$4 billion \parencite{AuditorGeneral2017} \\


\end{tabular}


%%% TABLE E.15 — Part (b): Spending + Social capital %%%

\centering
\footnotesize
\begin{tabular}{@{}
  >{\raggedright\arraybackslash}p{0.29\textwidth}
  >{\raggedright\arraybackslash}p{0.30\textwidth}
  >{\raggedright\arraybackslash}p{0.33\textwidth}
@{}}
\toprule
\textbf{Variable} &
\textbf{Chile (Maule)} &
\textbf{New Zealand (Canterbury)} \\
\midrule

\multicolumn{3}{@{}l}{\textit{\textbf{C. Public reconstruction spending}}} \\
\addlinespace[2pt]

Govt reconstruction commitment &
US\$8.4 billion &
NZ\$10.1 billion (incl.\ CERF + Natural Hazard Fund) \\

Timing of finance plan &
${\sim}$6 months (Aug 2010) &
${\sim}$3 months (May 2011 Budget) \\

Housing program scale &
220,000 units targeted &
${\sim}$70,000 managed repairs + Red Zone buyouts \\

\addlinespace[4pt]
\multicolumn{3}{@{}l}{\textit{\textbf{D. Informal institutions and social capital}}} \\
\addlinespace[2pt]

Pre-quake social trust &
Low \parencite{dussaillant2014trust} &
Moderate-to-high \\

Post-disaster civil disorder &
Moderate to high (post-disaster civil disorder and military curfew) &
Minimal \\

Community co-production &
Limited formal engagement &
Substantial (volunteers, entrepreneurs) \parencite{grube2018embedded} \\

\bottomrule
\end{tabular}

\vspace{4pt}
{\scriptsize \textit{Table~\ref{tab:quantcontrasts} concluded.} \textbf{Notes:} Insurance penetration (Chile): \parencite{AIRWorldwide2010}; \textcite{NguyenNoy2020} report residential coverage below 20\%. New~Zealand: near-universal, with roughly 80\% insured \parencite{NguyenNoy2020}. Reinsurance (NZ): \parencite{NHC2024}. CERA spending: \parencite{AuditorGeneral2017}. Red Zone buyouts: \parencite{Wilson2016}. Social capital: \parencite{dussaillant2014trust}. Settlement times: \parencite{Khakurel2021}. Over-cap transfers: \parencite{ICNZ2024}.}
\end{table}


% ============================================================
% E.3   REGULATORY AND LEGISLATIVE MILESTONE TIMELINE
% ============================================================
\subsection{Regulatory and Legislative Milestone Timeline}

Table~\ref{tab:regmilestones} provides a month-by-month comparison of the regulatory and legislative milestones in each country's disaster recovery. The contrast in institutional readiness is pronounced: New Zealand enacted emergency recovery legislation within 10~days of the September 2010 event, established CERA within 35~days of the February 2011 shock, and had gazetted technical land categories, building consent fast-tracking, and a public--private infrastructure alliance within 18~months. In contrast, Chile's response operated through pre-existing ministerial structures \textit{without} creating a dedicated recovery entity, enacting special legislative powers, or establishing formal regulatory fast-tracking for reconstruction.

%%% TABLE E.16 — Part (a): Month 0 through 8--12 %%%
\begin{table}[!htbp]
\centering
\caption{Regulatory and legislative milestone timeline (months from first event)}
\label{tab:regmilestones}
\footnotesize
\begin{tabular}{@{}
  >{\raggedright\arraybackslash}p{0.06\textwidth}
  >{\raggedright\arraybackslash}p{0.12\textwidth}
  >{\raggedright\arraybackslash}p{0.34\textwidth}
  >{\raggedright\arraybackslash}p{0.38\textwidth}
@{}}
\toprule
\textbf{Mo.} &
\textbf{Date} &
\textbf{Chile} &
\textbf{New Zealand} \\
\midrule

{0} & Feb--Sep 2010 &
State of Catastrophe (Feb~28). ONEMI response. No new legislation. &
Darfield earthquake (Sep~4). State of Emergency. Canterbury EQ Response and Recovery Act 2010 enacted (Sep~14). \\
\addlinespace

{1--2} & Oct--Nov 2010 &
Damage assessment ongoing. No new legislation. &
Canterbury EQ Recovery Commission established. EQC Catastrophe Response Plan operational. Fletcher Construction appointed. \\
\addlinespace

{3--6} & Dec 2010--Feb 2011 &
Plan published Aug 2010; no further legislative action. &
Feb~22 Christchurch earthquake (Mw~6.3, 185 deaths). National state of emergency declared. \\
\addlinespace

{6--7} & Mar--Apr 2011 &
Reconstruction through MINVU and existing ministers. No dedicated authority. &
CERA established 29 Mar 2011 (Order in Council). Canterbury EQ Recovery Act 2011 enacted (Apr). Powers to suspend/amend regulations. \\
\addlinespace

{8--12} & May--Sep 2011 &
Housing subsidies via SERVIUs. No regulatory fast-tracking. &
Red Zone announced (Jun). Technical categories (TC1--TC3) gazetted. Consent fast-tracking. 8,060 properties offered buyouts. \\

\end{tabular}

\centering
\footnotesize
\begin{tabular}{@{}
  >{\raggedright\arraybackslash}p{0.06\textwidth}
  >{\raggedright\arraybackslash}p{0.12\textwidth}
  >{\raggedright\arraybackslash}p{0.34\textwidth}
  >{\raggedright\arraybackslash}p{0.38\textwidth}
@{}}
\toprule
\textbf{Mo.} &
\textbf{Date} &
\textbf{Chile} &
\textbf{New Zealand} \\
\midrule

{12--18} & Oct 2011--Mar 2012 &
54\% of housing underway or complete by Oct 2012. No new institutional changes. &
EQC disputes-resolution manager appointed (Mar~2012). CCDU established within CERA (Apr~2012) for CBD recovery plan. \\
\addlinespace

{18--24} & Apr--Sep 2012 &
Ongoing centralized delivery. No public--private alliance formed. &
Canterbury Recovery Strategy approved (May~2012). SCIRT alliance operational. Budget ${\sim}$NZ\$2.5bn (actual: NZ\$2.2bn). \\
\addlinespace

{24--36} & Oct 2012--Sep 2013 &
54\% housing complete; 30\% under construction (Oct 2012). &
Land Use Recovery Plan approved (2013). Resilience integrated into district plans. Over 60 Orders in Council amending legislation. \\
\addlinespace

{36--60} & 2013--2016 &
Continued central govt transfers. No comparable construction surge. &
Peak construction 2015--2016. CERA spent NZ\$4bn. CERA disestablished 18~Apr 2016; functions transferred. \\

\bottomrule
\end{tabular}

\vspace{2pt}
{\scriptsize \textit{Table~\ref{tab:regmilestones} concluded. Sources:} Chile: \parencite{Comerio2013,Platt2019}. New Zealand: \parencite{CERRA2010,CERAA2011,AuditorGeneral2017,Wilson2016}.}
\end{table}


% ============================================================
% E.4   INTEGRATION OF MACRO-INSTITUTIONAL INDICATORS
% ============================================================
\subsection{Integration of Macro-Institutional Indicators (Table~\ref{tab:fiscal_inst}) with Disaster-Specific Evidence}

Table~\ref{tab:fiscal_inst} in the main text reports pre-quake national-level fiscal and institutional indicators for both countries. As noted in Section~5, these national-level scores serve as contextual descriptors rather than causal inputs; sub-national institutional quality indices are not available for the affected regions. Nevertheless, interpreting Table~\ref{tab:fiscal_inst} alongside the disaster-specific institutional evidence documented above yields several interpretively useful contrasts.

\textbf{Fiscal capacity and the insurance--debt inversion:} The most striking contrast in Table~\ref{tab:fiscal_inst} is between government debt (Chile 11.1\% of GDP; New Zealand 62.6\%) and the insurance architecture documented in Table~\ref{tab:quantcontrasts}. Chile entered the disaster with greater fiscal space, yet Maule's GDP trajectory shows no positive overshoot. New Zealand, despite higher sovereign indebtedness, channeled reconstruction capital primarily through its pre-existing insurance system rather than fiscal expansion. The magnitudes involved are telling: NZ's total insured cost reached NZ\$31~billion, of which more than NZ\$5~billion flowed from international reinsurers (Table~\ref{tab:quantcontrasts}), an exogenous capital injection that required no fiscal outlay. Chile's total insured cost was only US\$4--7~billion, with residential penetration below 20\%, leaving a US\$8.4~billion fiscal commitment as the primary reconstruction channel. Accordingly, Chile's change in public spending in the quake year (+10.44\%) exceeded New Zealand's (+6.24\%), directly reflecting this heavier reliance on the fiscal channel. The chronology in Table~\ref{tab:chronology} reinforces the point: NZ's EQC activated its claims process within days and appointed Fletcher Construction within weeks, whereas Chile's Plan de Reconstrucci{\'o}n was published six months after the earthquake. This contrast suggests that what matters for recovery dynamics is not necessarily raw fiscal capacity \textit{per se}, but the institutional and financial architecture through which reconstruction capital is mobilized. This observation does not imply that fiscal capacity is irrelevant, but rather that the \textit{speed and institutional channel} of financing may matter more for regional output than aggregate fiscal capacity.

\textbf{Legal system, property rights, and regulatory quality:} New Zealand scores higher on Legal System \& Property Rights (8.69 vs.\ 6.80) and Regulation (8.92 vs.\ 7.55). These national-level scores are consistent with the disaster-specific institutional evidence documented in Tables~\ref{tab:chronology} and~\ref{tab:regmilestones}: New Zealand enacted the Canterbury Earthquake Response and Recovery Act within 10~days of the September 2010 event, established CERA with statutory override powers within 35~days of the February 2011 shock, gazetted technical land categories and building consent fast-tracking within 12~months, and had the SCIRT public--private alliance and Land Use Recovery Plan operational within 24~months. Over CERA's lifespan, more than 60 Orders in Council were made to amend existing legislation for recovery purposes (Table~\ref{tab:regmilestones}). Chile, by contrast, made no comparable regulatory adaptations: reconstruction was managed through existing MINVU and MOP structures without a dedicated recovery entity, special legislative powers, or consent fast-tracking. The higher scores on Legal System \& Property Rights also align with the insurance channel: a well-functioning property rights regime is a \textit{precondition} for the complex contractual architecture of catastrophe insurance, reinsurance treaties, and the claims adjudication process through which more than 470,000 EQC and 168,000 private insurer claims were processed.

\textbf{Size of Government:} Chile's higher score (7.91 vs.\ 4.91) indicates a smaller government footprint relative to GDP (lower taxation, less government consumption, and fewer transfers). This is consistent with Chile's reliance on centralized but \textit{ad hoc} fiscal programs for reconstruction: the US\$8.4~billion commitment was channeled through existing ministry structures with no new institutional entities, whereas New Zealand's larger public sector provided the institutional infrastructure through which recovery was coordinated. Specifically, EQC (a Crown entity established in 1945 with NZ\$6.1~billion in accumulated reserves by 2010), CERA (a new government department with NZ\$4~billion in spending over five years), and the SCIRT alliance (a public--private arrangement coordinating three government entities and five construction firms on a NZ\$2.2~billion programme) were all creatures of New Zealand's comparatively larger state apparatus. The contrast illustrates that the relevant dimension is not government size in the abstract but rather the pre-existing institutional architecture that can be activated and adapted in a crisis. Chile's smaller government had correspondingly less pre-existing institutional capacity for disaster coordination (no public insurance scheme, no dedicated recovery authority) and the six-month lag to the Plan de Reconstrucci{\'o}n (Table~\ref{tab:chronology}) reflects the time required to build such capacity from scratch within existing ministerial structures.

To conclude, interpreting Table~\ref{tab:fiscal_inst} alongside the disaster-specific evidence in Tables~\ref{tab:chronology} to \ref{tab:regmilestones} reinforces the core interpretive point of Section~5: the divergent recovery trajectories are most consistently explained not by differential fiscal capacity (which actually favored Chile, at 11\% vs.\ 63\% debt-to-GDP) nor by broad macroeconomic stability, but by the pre-existing institutional architecture for disaster-specific financial mobilization, regulatory adaptation, and coordinated implementation. Chile's superior fiscal position did not translate into a comparable output response because the institutional channel through which spending was deployed (i.e., centralized ministerial programs without a dedicated recovery authority, or insurance-based external capital) was slower and less regionally targeted. The six-month lag to the Plan de Reconstrucci{\'o}n, the absence of regulatory override powers, and the limited community co-production all contrast with New Zealand's 35-day CERA establishment, over 60 Orders in Council adapting legislation, and the SCIRT public--private alliance.
 from your main .tex file.
%
% REQUIRED in your main preamble (most likely already loaded):
%   \usepackage{booktabs}
%   \usepackage{array}
%   \usepackage{hyperref}
% ============================================================

\renewcommand{\thetable}{E.\arabic{table}}
\setcounter{table}{13}
\setcounter{subsection}{0}

\section{Institutional context and comparison: quantitative summary}
\addcontentsline{toc}{section}{Appendix E: Institutional Context}

This appendix documents the factual institutional context underlying the interpretive discussion in Section~5. It provides: (i)~a side-by-side chronology of the formal institutional responses, (ii)~a summary table of quantitative institutional contrasts, (iii)~a regulatory milestone timeline, and (iv)~an interpretive integration of the Table~\ref{tab:fiscal_inst} with the disaster-specific evidence. Importantly, these are descriptive inputs, not causal estimates; they serve to ground the mechanism discussion in observable facts.


% ============================================================
% E.1   COMPARATIVE INSTITUTIONAL CHRONOLOGY
% ============================================================
\subsection{Comparative Institutional Chronology}

Table~\ref{tab:chronology} presents a side-by-side timeline of key institutional actions.

%%% TABLE E.14 — Part (a) %%%
\begin{table}[!htbp]
\centering
\caption{Side-by-side institutional response chronology}
\label{tab:chronology}
\footnotesize
\begin{tabular}{@{}
  >{\raggedright\arraybackslash}p{0.12\textwidth}
  >{\raggedright\arraybackslash}p{0.40\textwidth}
  >{\raggedright\arraybackslash}p{0.40\textwidth}
@{}}
\toprule
\textbf{Timeframe} &
\textbf{Chile (Maule, Feb 2010)} &
\textbf{New Zealand (Canterbury, Sep 2010--Feb 2011)} \\
\midrule

Days 1--7} &
State of Catastrophe declared in Maule and Biob{\'{\i}}o (Feb~28). Military curfew imposed. Approx.\ 14,000 troops deployed \parencite{CRS2010}. ONEMI coordinated initial response; later criticized for tsunami warning failures. &
Sep 2010 (Darfield): State of Emergency declared. Civil Defence managed initial response. No significant civil disorder. EQC activated Catastrophe Response Plan; claims lodging opened within days. \\
\addlinespace

{Weeks 2--8} &
Presidential transition (Bachelet to Pi{\~n}era, Mar~11) complicated coordination. National reconstruction coordinator appointed at MINVU. Initial damage surveys began. No centralized claims system for the approx.\ 10--20\% of residentially insured properties. &
EQC appointed Fletcher Construction to project-manage residential repairs (approx.\ 50,000 homes). Assessment teams deployed. Small claims (under-cap) assessed within months \parencite{Khakurel2021}. Private insurers processing in parallel. \\
\addlinespace

{Months 2--6} &
Plan de Reconstrucci{\'o}n published by MINVU (Aug~2010), six months post-quake. Targeted 220,000 housing units over four years. Finance plan adopted. No dedicated recovery authority created. &
Feb 2011 Christchurch earthquake (Mw~6.3, 185 deaths). CERA established 29~Mar 2011---35 days after the Feb event. Wide-ranging powers including authority to suspend laws. Canterbury Earthquake Recovery Act passed (Apr~2011). \\

\bottomrule
\end{tabular}

\vspace{2pt}
{\scriptsize\textit{Continued in next panel.}}
\end{table}

%%% TABLE E.14 — Part (b) %%%
\begin{table}[!htbp]
\centering
\footnotesize
\begin{tabular}{@{}
  >{\raggedright\arraybackslash}p{0.12\textwidth}
  >{\raggedright\arraybackslash}p{0.40\textwidth}
  >{\raggedright\arraybackslash}p{0.40\textwidth}
@{}}
\toprule
\textbf{Timeframe} &
\textbf{Chile (Maule, Feb 2010)} &
\textbf{New Zealand (Canterbury, Sep 2010--Feb 2011)} \\
\midrule

{Months 6--12} &
Reconstruction managed through existing ministries (MINVU, MOP). No new entity created. Some temporary shelters operational. &
Red Zone announced; purchase offers to approx.\ 8,060 properties. Technical categories (TC1--TC3) established. Building consent fast-tracking initiated. Fletcher EQR expanded to approx.\ 70,000 homes. \\
\addlinespace

{Year 2} &
Centralized ministerial direction continued. No significant regulatory changes. Community participation limited \parencite{CardenasJiron2012}. By Oct~2012: 54\% housing complete; 30\% under construction \parencite{Platt2019}. &
Canterbury Recovery Strategy approved (May~2012). SCIRT alliance formed (CERA + CCC + NZTA + 5 firms). Land Use Recovery Plan approved (2013). Community groups active. Peak construction began. \\
\addlinespace

{Years 3--5} &
Spending continued through central transfers. No comparable construction sector surge relative to counterfactual (Section~\ref{sec:sectoral_impact_analysis}). &
CERA spent NZ\$4bn over 5~years. Construction share peaked ${\sim}$10\% in 2016. CERA disestablished 18~Apr 2016; functions transferred to DPMC, LINZ, MBIE, CCC, Regenerate Christchurch, {\=O}takaro Ltd. \\

\bottomrule
\end{tabular}

\vspace{2pt}
{\scriptsize \textit{Table~\ref{tab:chronology} concluded. Sources:} Chile: \parencite{eclac2010overview,Platt2019,Comerio2013,CRS2010}. New Zealand: \parencite{AuditorGeneral2017,CERAA2011,ICNZ2024,eqc2017annual,grube2018embedded,Wilson2016}.}
\end{table}


% ============================================================
% E.2   QUANTITATIVE INSTITUTIONAL CONTRASTS
% ============================================================
\subsection{Quantitative Institutional Contrasts}

Table~\ref{tab:quantcontrasts} summarizes the key quantitative differences in institutional architecture available for post-disaster recovery.

%%% TABLE E.15 — Part (a): Insurance + Governance %%%
\begin{table}[!htbp]
\centering
\caption{Key quantitative institutional contrasts: Chile vs.\ New Zealand}
\label{tab:quantcontrasts}
\footnotesize
\begin{tabular}{@{}
  >{\raggedright\arraybackslash}p{0.29\textwidth}
  >{\raggedright\arraybackslash}p{0.30\textwidth}
  >{\raggedright\arraybackslash}p{0.33\textwidth}
@{}}
\toprule
\textbf{Variable} &
\textbf{Chile (Maule)} &
\textbf{New Zealand (Canterbury)} \\
\midrule

\multicolumn{3}{@{}l}{\textit{\textbf{A. Insurance architecture}}} \\
\addlinespace[2pt]

Residential EQ insurance penetration &
${\sim}$10--20\% &
Near-universal (EQC automatic with fire ins.) \\

Commercial insurance penetration &
${\sim}$60\% &
High (bundled) \\

Pre-existing public insurance scheme &
None &
EQC (est.\ 1945) \\

Pre-quake accumulated fund &
n/a &
NZ\$6.1 billion \parencite{eqc2017annual} \\

Total insured cost &
US\$4--7 billion &
NZ\$31 billion \parencite{ICNZ2024} \\

EQC claims filed &
n/a &
$>$470,000 \\

Intl.\ reinsurance recovered &
n/a &
$>$NZ\$5 billion \parencite{NHC2024} \\

Median settlement time (EQC) &
No comparable public data &
Small: 138--191\,d; Med: 712--890\,d; Large: 1,226--1,392\,d \parencite{Khakurel2021} \\

Crown guarantee for shortfall &
None &
Yes (EQC Act s.\ 16) \\

\addlinespace[4pt]
\multicolumn{3}{@{}l}{\textit{\textbf{B. Governance and recovery institutions}}} \\
\addlinespace[2pt]

Dedicated recovery authority &
None created &
CERA (est.\ 29 Mar 2011; disest.\ 18 Apr 2016) \\

Time to recovery authority &
--- &
35 days \\

Special legislative powers &
None enacted &
Canterbury EQ Recovery Act 2011 \\

Regulatory fast-tracking &
No significant changes &
Yes (technical categories, consent fast-tracking) \\

Public--private rebuild alliance &
No &
SCIRT (CCC + CERA + NZTA + 5 firms) \\

Recovery authority spending (5~yrs) &
--- &
NZ\$4 billion \parencite{AuditorGeneral2017} \\


\end{tabular}


%%% TABLE E.15 — Part (b): Spending + Social capital %%%

\centering
\footnotesize
\begin{tabular}{@{}
  >{\raggedright\arraybackslash}p{0.29\textwidth}
  >{\raggedright\arraybackslash}p{0.30\textwidth}
  >{\raggedright\arraybackslash}p{0.33\textwidth}
@{}}
\toprule
\textbf{Variable} &
\textbf{Chile (Maule)} &
\textbf{New Zealand (Canterbury)} \\
\midrule

\multicolumn{3}{@{}l}{\textit{\textbf{C. Public reconstruction spending}}} \\
\addlinespace[2pt]

Govt reconstruction commitment &
US\$8.4 billion &
NZ\$10.1 billion (incl.\ CERF + Natural Hazard Fund) \\

Timing of finance plan &
${\sim}$6 months (Aug 2010) &
${\sim}$3 months (May 2011 Budget) \\

Housing program scale &
220,000 units targeted &
${\sim}$70,000 managed repairs + Red Zone buyouts \\

\addlinespace[4pt]
\multicolumn{3}{@{}l}{\textit{\textbf{D. Informal institutions and social capital}}} \\
\addlinespace[2pt]

Pre-quake social trust &
Low \parencite{dussaillant2014trust} &
Moderate-to-high \\

Post-disaster civil disorder &
Moderate to high (post-disaster civil disorder and military curfew) &
Minimal \\

Community co-production &
Limited formal engagement &
Substantial (volunteers, entrepreneurs) \parencite{grube2018embedded} \\

\bottomrule
\end{tabular}

\vspace{4pt}
{\scriptsize \textit{Table~\ref{tab:quantcontrasts} concluded.} \textbf{Notes:} Insurance penetration (Chile): \parencite{AIRWorldwide2010}; \textcite{NguyenNoy2020} report residential coverage below 20\%. New~Zealand: near-universal, with roughly 80\% insured \parencite{NguyenNoy2020}. Reinsurance (NZ): \parencite{NHC2024}. CERA spending: \parencite{AuditorGeneral2017}. Red Zone buyouts: \parencite{Wilson2016}. Social capital: \parencite{dussaillant2014trust}. Settlement times: \parencite{Khakurel2021}. Over-cap transfers: \parencite{ICNZ2024}.}
\end{table}


% ============================================================
% E.3   REGULATORY AND LEGISLATIVE MILESTONE TIMELINE
% ============================================================
\subsection{Regulatory and Legislative Milestone Timeline}

Table~\ref{tab:regmilestones} provides a month-by-month comparison of the regulatory and legislative milestones in each country's disaster recovery. The contrast in institutional readiness is pronounced: New Zealand enacted emergency recovery legislation within 10~days of the September 2010 event, established CERA within 35~days of the February 2011 shock, and had gazetted technical land categories, building consent fast-tracking, and a public--private infrastructure alliance within 18~months. In contrast, Chile's response operated through pre-existing ministerial structures \textit{without} creating a dedicated recovery entity, enacting special legislative powers, or establishing formal regulatory fast-tracking for reconstruction.

%%% TABLE E.16 — Part (a): Month 0 through 8--12 %%%
\begin{table}[!htbp]
\centering
\caption{Regulatory and legislative milestone timeline (months from first event)}
\label{tab:regmilestones}
\footnotesize
\begin{tabular}{@{}
  >{\raggedright\arraybackslash}p{0.06\textwidth}
  >{\raggedright\arraybackslash}p{0.12\textwidth}
  >{\raggedright\arraybackslash}p{0.34\textwidth}
  >{\raggedright\arraybackslash}p{0.38\textwidth}
@{}}
\toprule
\textbf{Mo.} &
\textbf{Date} &
\textbf{Chile} &
\textbf{New Zealand} \\
\midrule

{0} & Feb--Sep 2010 &
State of Catastrophe (Feb~28). ONEMI response. No new legislation. &
Darfield earthquake (Sep~4). State of Emergency. Canterbury EQ Response and Recovery Act 2010 enacted (Sep~14). \\
\addlinespace

{1--2} & Oct--Nov 2010 &
Damage assessment ongoing. No new legislation. &
Canterbury EQ Recovery Commission established. EQC Catastrophe Response Plan operational. Fletcher Construction appointed. \\
\addlinespace

{3--6} & Dec 2010--Feb 2011 &
Plan published Aug 2010; no further legislative action. &
Feb~22 Christchurch earthquake (Mw~6.3, 185 deaths). National state of emergency declared. \\
\addlinespace

{6--7} & Mar--Apr 2011 &
Reconstruction through MINVU and existing ministers. No dedicated authority. &
CERA established 29 Mar 2011 (Order in Council). Canterbury EQ Recovery Act 2011 enacted (Apr). Powers to suspend/amend regulations. \\
\addlinespace

{8--12} & May--Sep 2011 &
Housing subsidies via SERVIUs. No regulatory fast-tracking. &
Red Zone announced (Jun). Technical categories (TC1--TC3) gazetted. Consent fast-tracking. 8,060 properties offered buyouts. \\

\end{tabular}

\centering
\footnotesize
\begin{tabular}{@{}
  >{\raggedright\arraybackslash}p{0.06\textwidth}
  >{\raggedright\arraybackslash}p{0.12\textwidth}
  >{\raggedright\arraybackslash}p{0.34\textwidth}
  >{\raggedright\arraybackslash}p{0.38\textwidth}
@{}}
\toprule
\textbf{Mo.} &
\textbf{Date} &
\textbf{Chile} &
\textbf{New Zealand} \\
\midrule

{12--18} & Oct 2011--Mar 2012 &
54\% of housing underway or complete by Oct 2012. No new institutional changes. &
EQC disputes-resolution manager appointed (Mar~2012). CCDU established within CERA (Apr~2012) for CBD recovery plan. \\
\addlinespace

{18--24} & Apr--Sep 2012 &
Ongoing centralized delivery. No public--private alliance formed. &
Canterbury Recovery Strategy approved (May~2012). SCIRT alliance operational. Budget ${\sim}$NZ\$2.5bn (actual: NZ\$2.2bn). \\
\addlinespace

{24--36} & Oct 2012--Sep 2013 &
54\% housing complete; 30\% under construction (Oct 2012). &
Land Use Recovery Plan approved (2013). Resilience integrated into district plans. Over 60 Orders in Council amending legislation. \\
\addlinespace

{36--60} & 2013--2016 &
Continued central govt transfers. No comparable construction surge. &
Peak construction 2015--2016. CERA spent NZ\$4bn. CERA disestablished 18~Apr 2016; functions transferred. \\

\bottomrule
\end{tabular}

\vspace{2pt}
{\scriptsize \textit{Table~\ref{tab:regmilestones} concluded. Sources:} Chile: \parencite{Comerio2013,Platt2019}. New Zealand: \parencite{CERRA2010,CERAA2011,AuditorGeneral2017,Wilson2016}.}
\end{table}


% ============================================================
% E.4   INTEGRATION OF MACRO-INSTITUTIONAL INDICATORS
% ============================================================
\subsection{Integration of Macro-Institutional Indicators (Table~\ref{tab:fiscal_inst}) with Disaster-Specific Evidence}

Table~\ref{tab:fiscal_inst} in the main text reports pre-quake national-level fiscal and institutional indicators for both countries. As noted in Section~5, these national-level scores serve as contextual descriptors rather than causal inputs; sub-national institutional quality indices are not available for the affected regions. Nevertheless, interpreting Table~\ref{tab:fiscal_inst} alongside the disaster-specific institutional evidence documented above yields several interpretively useful contrasts.

\textbf{Fiscal capacity and the insurance--debt inversion:} The most striking contrast in Table~\ref{tab:fiscal_inst} is between government debt (Chile 11.1\% of GDP; New Zealand 62.6\%) and the insurance architecture documented in Table~\ref{tab:quantcontrasts}. Chile entered the disaster with greater fiscal space, yet Maule's GDP trajectory shows no positive overshoot. New Zealand, despite higher sovereign indebtedness, channeled reconstruction capital primarily through its pre-existing insurance system rather than fiscal expansion. The magnitudes involved are telling: NZ's total insured cost reached NZ\$31~billion, of which more than NZ\$5~billion flowed from international reinsurers (Table~\ref{tab:quantcontrasts}), an exogenous capital injection that required no fiscal outlay. Chile's total insured cost was only US\$4--7~billion, with residential penetration below 20\%, leaving a US\$8.4~billion fiscal commitment as the primary reconstruction channel. Accordingly, Chile's change in public spending in the quake year (+10.44\%) exceeded New Zealand's (+6.24\%), directly reflecting this heavier reliance on the fiscal channel. The chronology in Table~\ref{tab:chronology} reinforces the point: NZ's EQC activated its claims process within days and appointed Fletcher Construction within weeks, whereas Chile's Plan de Reconstrucci{\'o}n was published six months after the earthquake. This contrast suggests that what matters for recovery dynamics is not necessarily raw fiscal capacity \textit{per se}, but the institutional and financial architecture through which reconstruction capital is mobilized. This observation does not imply that fiscal capacity is irrelevant, but rather that the \textit{speed and institutional channel} of financing may matter more for regional output than aggregate fiscal capacity.

\textbf{Legal system, property rights, and regulatory quality:} New Zealand scores higher on Legal System \& Property Rights (8.69 vs.\ 6.80) and Regulation (8.92 vs.\ 7.55). These national-level scores are consistent with the disaster-specific institutional evidence documented in Tables~\ref{tab:chronology} and~\ref{tab:regmilestones}: New Zealand enacted the Canterbury Earthquake Response and Recovery Act within 10~days of the September 2010 event, established CERA with statutory override powers within 35~days of the February 2011 shock, gazetted technical land categories and building consent fast-tracking within 12~months, and had the SCIRT public--private alliance and Land Use Recovery Plan operational within 24~months. Over CERA's lifespan, more than 60 Orders in Council were made to amend existing legislation for recovery purposes (Table~\ref{tab:regmilestones}). Chile, by contrast, made no comparable regulatory adaptations: reconstruction was managed through existing MINVU and MOP structures without a dedicated recovery entity, special legislative powers, or consent fast-tracking. The higher scores on Legal System \& Property Rights also align with the insurance channel: a well-functioning property rights regime is a \textit{precondition} for the complex contractual architecture of catastrophe insurance, reinsurance treaties, and the claims adjudication process through which more than 470,000 EQC and 168,000 private insurer claims were processed.

\textbf{Size of Government:} Chile's higher score (7.91 vs.\ 4.91) indicates a smaller government footprint relative to GDP (lower taxation, less government consumption, and fewer transfers). This is consistent with Chile's reliance on centralized but \textit{ad hoc} fiscal programs for reconstruction: the US\$8.4~billion commitment was channeled through existing ministry structures with no new institutional entities, whereas New Zealand's larger public sector provided the institutional infrastructure through which recovery was coordinated. Specifically, EQC (a Crown entity established in 1945 with NZ\$6.1~billion in accumulated reserves by 2010), CERA (a new government department with NZ\$4~billion in spending over five years), and the SCIRT alliance (a public--private arrangement coordinating three government entities and five construction firms on a NZ\$2.2~billion programme) were all creatures of New Zealand's comparatively larger state apparatus. The contrast illustrates that the relevant dimension is not government size in the abstract but rather the pre-existing institutional architecture that can be activated and adapted in a crisis. Chile's smaller government had correspondingly less pre-existing institutional capacity for disaster coordination (no public insurance scheme, no dedicated recovery authority) and the six-month lag to the Plan de Reconstrucci{\'o}n (Table~\ref{tab:chronology}) reflects the time required to build such capacity from scratch within existing ministerial structures.

To conclude, interpreting Table~\ref{tab:fiscal_inst} alongside the disaster-specific evidence in Tables~\ref{tab:chronology} to \ref{tab:regmilestones} reinforces the core interpretive point of Section~5: the divergent recovery trajectories are most consistently explained not by differential fiscal capacity (which actually favored Chile, at 11\% vs.\ 63\% debt-to-GDP) nor by broad macroeconomic stability, but by the pre-existing institutional architecture for disaster-specific financial mobilization, regulatory adaptation, and coordinated implementation. Chile's superior fiscal position did not translate into a comparable output response because the institutional channel through which spending was deployed (i.e., centralized ministerial programs without a dedicated recovery authority, or insurance-based external capital) was slower and less regionally targeted. The six-month lag to the Plan de Reconstrucci{\'o}n, the absence of regulatory override powers, and the limited community co-production all contrast with New Zealand's 35-day CERA establishment, over 60 Orders in Council adapting legislation, and the SCIRT public--private alliance.
 from your main .tex file.
%
% REQUIRED in your main preamble (most likely already loaded):
%   \usepackage{booktabs}
%   \usepackage{array}
%   \usepackage{hyperref}
% ============================================================

\renewcommand{\thetable}{E.\arabic{table}}
\setcounter{table}{13}
\setcounter{subsection}{0}

\section{Institutional context and comparison: quantitative summary}
\addcontentsline{toc}{section}{Appendix E: Institutional Context}

This appendix documents the factual institutional context underlying the interpretive discussion in Section~5. It provides: (i)~a side-by-side chronology of the formal institutional responses, (ii)~a summary table of quantitative institutional contrasts, (iii)~a regulatory milestone timeline, and (iv)~an interpretive integration of the Table~\ref{tab:fiscal_inst} with the disaster-specific evidence. Importantly, these are descriptive inputs, not causal estimates; they serve to ground the mechanism discussion in observable facts.


% ============================================================
% E.1   COMPARATIVE INSTITUTIONAL CHRONOLOGY
% ============================================================
\subsection{Comparative Institutional Chronology}

Table~\ref{tab:chronology} presents a side-by-side timeline of key institutional actions.

%%% TABLE E.14 — Part (a) %%%
\begin{table}[!htbp]
\centering
\caption{Side-by-side institutional response chronology}
\label{tab:chronology}
\footnotesize
\begin{tabular}{@{}
  >{\raggedright\arraybackslash}p{0.12\textwidth}
  >{\raggedright\arraybackslash}p{0.40\textwidth}
  >{\raggedright\arraybackslash}p{0.40\textwidth}
@{}}
\toprule
\textbf{Timeframe} &
\textbf{Chile (Maule, Feb 2010)} &
\textbf{New Zealand (Canterbury, Sep 2010--Feb 2011)} \\
\midrule

Days 1--7} &
State of Catastrophe declared in Maule and Biob{\'{\i}}o (Feb~28). Military curfew imposed. Approx.\ 14,000 troops deployed \parencite{CRS2010}. ONEMI coordinated initial response; later criticized for tsunami warning failures. &
Sep 2010 (Darfield): State of Emergency declared. Civil Defence managed initial response. No significant civil disorder. EQC activated Catastrophe Response Plan; claims lodging opened within days. \\
\addlinespace

{Weeks 2--8} &
Presidential transition (Bachelet to Pi{\~n}era, Mar~11) complicated coordination. National reconstruction coordinator appointed at MINVU. Initial damage surveys began. No centralized claims system for the approx.\ 10--20\% of residentially insured properties. &
EQC appointed Fletcher Construction to project-manage residential repairs (approx.\ 50,000 homes). Assessment teams deployed. Small claims (under-cap) assessed within months \parencite{Khakurel2021}. Private insurers processing in parallel. \\
\addlinespace

{Months 2--6} &
Plan de Reconstrucci{\'o}n published by MINVU (Aug~2010), six months post-quake. Targeted 220,000 housing units over four years. Finance plan adopted. No dedicated recovery authority created. &
Feb 2011 Christchurch earthquake (Mw~6.3, 185 deaths). CERA established 29~Mar 2011---35 days after the Feb event. Wide-ranging powers including authority to suspend laws. Canterbury Earthquake Recovery Act passed (Apr~2011). \\

\bottomrule
\end{tabular}

\vspace{2pt}
{\scriptsize\textit{Continued in next panel.}}
\end{table}

%%% TABLE E.14 — Part (b) %%%
\begin{table}[!htbp]
\centering
\footnotesize
\begin{tabular}{@{}
  >{\raggedright\arraybackslash}p{0.12\textwidth}
  >{\raggedright\arraybackslash}p{0.40\textwidth}
  >{\raggedright\arraybackslash}p{0.40\textwidth}
@{}}
\toprule
\textbf{Timeframe} &
\textbf{Chile (Maule, Feb 2010)} &
\textbf{New Zealand (Canterbury, Sep 2010--Feb 2011)} \\
\midrule

{Months 6--12} &
Reconstruction managed through existing ministries (MINVU, MOP). No new entity created. Some temporary shelters operational. &
Red Zone announced; purchase offers to approx.\ 8,060 properties. Technical categories (TC1--TC3) established. Building consent fast-tracking initiated. Fletcher EQR expanded to approx.\ 70,000 homes. \\
\addlinespace

{Year 2} &
Centralized ministerial direction continued. No significant regulatory changes. Community participation limited \parencite{CardenasJiron2012}. By Oct~2012: 54\% housing complete; 30\% under construction \parencite{Platt2019}. &
Canterbury Recovery Strategy approved (May~2012). SCIRT alliance formed (CERA + CCC + NZTA + 5 firms). Land Use Recovery Plan approved (2013). Community groups active. Peak construction began. \\
\addlinespace

{Years 3--5} &
Spending continued through central transfers. No comparable construction sector surge relative to counterfactual (Section~\ref{sec:sectoral_impact_analysis}). &
CERA spent NZ\$4bn over 5~years. Construction share peaked ${\sim}$10\% in 2016. CERA disestablished 18~Apr 2016; functions transferred to DPMC, LINZ, MBIE, CCC, Regenerate Christchurch, {\=O}takaro Ltd. \\

\bottomrule
\end{tabular}

\vspace{2pt}
{\scriptsize \textit{Table~\ref{tab:chronology} concluded. Sources:} Chile: \parencite{eclac2010overview,Platt2019,Comerio2013,CRS2010}. New Zealand: \parencite{AuditorGeneral2017,CERAA2011,ICNZ2024,eqc2017annual,grube2018embedded,Wilson2016}.}
\end{table}


% ============================================================
% E.2   QUANTITATIVE INSTITUTIONAL CONTRASTS
% ============================================================
\subsection{Quantitative Institutional Contrasts}

Table~\ref{tab:quantcontrasts} summarizes the key quantitative differences in institutional architecture available for post-disaster recovery.

%%% TABLE E.15 — Part (a): Insurance + Governance %%%
\begin{table}[!htbp]
\centering
\caption{Key quantitative institutional contrasts: Chile vs.\ New Zealand}
\label{tab:quantcontrasts}
\footnotesize
\begin{tabular}{@{}
  >{\raggedright\arraybackslash}p{0.29\textwidth}
  >{\raggedright\arraybackslash}p{0.30\textwidth}
  >{\raggedright\arraybackslash}p{0.33\textwidth}
@{}}
\toprule
\textbf{Variable} &
\textbf{Chile (Maule)} &
\textbf{New Zealand (Canterbury)} \\
\midrule

\multicolumn{3}{@{}l}{\textit{\textbf{A. Insurance architecture}}} \\
\addlinespace[2pt]

Residential EQ insurance penetration &
${\sim}$10--20\% &
Near-universal (EQC automatic with fire ins.) \\

Commercial insurance penetration &
${\sim}$60\% &
High (bundled) \\

Pre-existing public insurance scheme &
None &
EQC (est.\ 1945) \\

Pre-quake accumulated fund &
n/a &
NZ\$6.1 billion \parencite{eqc2017annual} \\

Total insured cost &
US\$4--7 billion &
NZ\$31 billion \parencite{ICNZ2024} \\

EQC claims filed &
n/a &
$>$470,000 \\

Intl.\ reinsurance recovered &
n/a &
$>$NZ\$5 billion \parencite{NHC2024} \\

Median settlement time (EQC) &
No comparable public data &
Small: 138--191\,d; Med: 712--890\,d; Large: 1,226--1,392\,d \parencite{Khakurel2021} \\

Crown guarantee for shortfall &
None &
Yes (EQC Act s.\ 16) \\

\addlinespace[4pt]
\multicolumn{3}{@{}l}{\textit{\textbf{B. Governance and recovery institutions}}} \\
\addlinespace[2pt]

Dedicated recovery authority &
None created &
CERA (est.\ 29 Mar 2011; disest.\ 18 Apr 2016) \\

Time to recovery authority &
--- &
35 days \\

Special legislative powers &
None enacted &
Canterbury EQ Recovery Act 2011 \\

Regulatory fast-tracking &
No significant changes &
Yes (technical categories, consent fast-tracking) \\

Public--private rebuild alliance &
No &
SCIRT (CCC + CERA + NZTA + 5 firms) \\

Recovery authority spending (5~yrs) &
--- &
NZ\$4 billion \parencite{AuditorGeneral2017} \\


\end{tabular}


%%% TABLE E.15 — Part (b): Spending + Social capital %%%

\centering
\footnotesize
\begin{tabular}{@{}
  >{\raggedright\arraybackslash}p{0.29\textwidth}
  >{\raggedright\arraybackslash}p{0.30\textwidth}
  >{\raggedright\arraybackslash}p{0.33\textwidth}
@{}}
\toprule
\textbf{Variable} &
\textbf{Chile (Maule)} &
\textbf{New Zealand (Canterbury)} \\
\midrule

\multicolumn{3}{@{}l}{\textit{\textbf{C. Public reconstruction spending}}} \\
\addlinespace[2pt]

Govt reconstruction commitment &
US\$8.4 billion &
NZ\$10.1 billion (incl.\ CERF + Natural Hazard Fund) \\

Timing of finance plan &
${\sim}$6 months (Aug 2010) &
${\sim}$3 months (May 2011 Budget) \\

Housing program scale &
220,000 units targeted &
${\sim}$70,000 managed repairs + Red Zone buyouts \\

\addlinespace[4pt]
\multicolumn{3}{@{}l}{\textit{\textbf{D. Informal institutions and social capital}}} \\
\addlinespace[2pt]

Pre-quake social trust &
Low \parencite{dussaillant2014trust} &
Moderate-to-high \\

Post-disaster civil disorder &
Moderate to high (post-disaster civil disorder and military curfew) &
Minimal \\

Community co-production &
Limited formal engagement &
Substantial (volunteers, entrepreneurs) \parencite{grube2018embedded} \\

\bottomrule
\end{tabular}

\vspace{4pt}
{\scriptsize \textit{Table~\ref{tab:quantcontrasts} concluded.} \textbf{Notes:} Insurance penetration (Chile): \parencite{AIRWorldwide2010}; \textcite{NguyenNoy2020} report residential coverage below 20\%. New~Zealand: near-universal, with roughly 80\% insured \parencite{NguyenNoy2020}. Reinsurance (NZ): \parencite{NHC2024}. CERA spending: \parencite{AuditorGeneral2017}. Red Zone buyouts: \parencite{Wilson2016}. Social capital: \parencite{dussaillant2014trust}. Settlement times: \parencite{Khakurel2021}. Over-cap transfers: \parencite{ICNZ2024}.}
\end{table}


% ============================================================
% E.3   REGULATORY AND LEGISLATIVE MILESTONE TIMELINE
% ============================================================
\subsection{Regulatory and Legislative Milestone Timeline}

Table~\ref{tab:regmilestones} provides a month-by-month comparison of the regulatory and legislative milestones in each country's disaster recovery. The contrast in institutional readiness is pronounced: New Zealand enacted emergency recovery legislation within 10~days of the September 2010 event, established CERA within 35~days of the February 2011 shock, and had gazetted technical land categories, building consent fast-tracking, and a public--private infrastructure alliance within 18~months. In contrast, Chile's response operated through pre-existing ministerial structures \textit{without} creating a dedicated recovery entity, enacting special legislative powers, or establishing formal regulatory fast-tracking for reconstruction.

%%% TABLE E.16 — Part (a): Month 0 through 8--12 %%%
\begin{table}[!htbp]
\centering
\caption{Regulatory and legislative milestone timeline (months from first event)}
\label{tab:regmilestones}
\footnotesize
\begin{tabular}{@{}
  >{\raggedright\arraybackslash}p{0.06\textwidth}
  >{\raggedright\arraybackslash}p{0.12\textwidth}
  >{\raggedright\arraybackslash}p{0.34\textwidth}
  >{\raggedright\arraybackslash}p{0.38\textwidth}
@{}}
\toprule
\textbf{Mo.} &
\textbf{Date} &
\textbf{Chile} &
\textbf{New Zealand} \\
\midrule

{0} & Feb--Sep 2010 &
State of Catastrophe (Feb~28). ONEMI response. No new legislation. &
Darfield earthquake (Sep~4). State of Emergency. Canterbury EQ Response and Recovery Act 2010 enacted (Sep~14). \\
\addlinespace

{1--2} & Oct--Nov 2010 &
Damage assessment ongoing. No new legislation. &
Canterbury EQ Recovery Commission established. EQC Catastrophe Response Plan operational. Fletcher Construction appointed. \\
\addlinespace

{3--6} & Dec 2010--Feb 2011 &
Plan published Aug 2010; no further legislative action. &
Feb~22 Christchurch earthquake (Mw~6.3, 185 deaths). National state of emergency declared. \\
\addlinespace

{6--7} & Mar--Apr 2011 &
Reconstruction through MINVU and existing ministers. No dedicated authority. &
CERA established 29 Mar 2011 (Order in Council). Canterbury EQ Recovery Act 2011 enacted (Apr). Powers to suspend/amend regulations. \\
\addlinespace

{8--12} & May--Sep 2011 &
Housing subsidies via SERVIUs. No regulatory fast-tracking. &
Red Zone announced (Jun). Technical categories (TC1--TC3) gazetted. Consent fast-tracking. 8,060 properties offered buyouts. \\

\end{tabular}

\centering
\footnotesize
\begin{tabular}{@{}
  >{\raggedright\arraybackslash}p{0.06\textwidth}
  >{\raggedright\arraybackslash}p{0.12\textwidth}
  >{\raggedright\arraybackslash}p{0.34\textwidth}
  >{\raggedright\arraybackslash}p{0.38\textwidth}
@{}}
\toprule
\textbf{Mo.} &
\textbf{Date} &
\textbf{Chile} &
\textbf{New Zealand} \\
\midrule

{12--18} & Oct 2011--Mar 2012 &
54\% of housing underway or complete by Oct 2012. No new institutional changes. &
EQC disputes-resolution manager appointed (Mar~2012). CCDU established within CERA (Apr~2012) for CBD recovery plan. \\
\addlinespace

{18--24} & Apr--Sep 2012 &
Ongoing centralized delivery. No public--private alliance formed. &
Canterbury Recovery Strategy approved (May~2012). SCIRT alliance operational. Budget ${\sim}$NZ\$2.5bn (actual: NZ\$2.2bn). \\
\addlinespace

{24--36} & Oct 2012--Sep 2013 &
54\% housing complete; 30\% under construction (Oct 2012). &
Land Use Recovery Plan approved (2013). Resilience integrated into district plans. Over 60 Orders in Council amending legislation. \\
\addlinespace

{36--60} & 2013--2016 &
Continued central govt transfers. No comparable construction surge. &
Peak construction 2015--2016. CERA spent NZ\$4bn. CERA disestablished 18~Apr 2016; functions transferred. \\

\bottomrule
\end{tabular}

\vspace{2pt}
{\scriptsize \textit{Table~\ref{tab:regmilestones} concluded. Sources:} Chile: \parencite{Comerio2013,Platt2019}. New Zealand: \parencite{CERRA2010,CERAA2011,AuditorGeneral2017,Wilson2016}.}
\end{table}


% ============================================================
% E.4   INTEGRATION OF MACRO-INSTITUTIONAL INDICATORS
% ============================================================
\subsection{Integration of Macro-Institutional Indicators (Table~\ref{tab:fiscal_inst}) with Disaster-Specific Evidence}

Table~\ref{tab:fiscal_inst} in the main text reports pre-quake national-level fiscal and institutional indicators for both countries. As noted in Section~5, these national-level scores serve as contextual descriptors rather than causal inputs; sub-national institutional quality indices are not available for the affected regions. Nevertheless, interpreting Table~\ref{tab:fiscal_inst} alongside the disaster-specific institutional evidence documented above yields several interpretively useful contrasts.

\textbf{Fiscal capacity and the insurance--debt inversion:} The most striking contrast in Table~\ref{tab:fiscal_inst} is between government debt (Chile 11.1\% of GDP; New Zealand 62.6\%) and the insurance architecture documented in Table~\ref{tab:quantcontrasts}. Chile entered the disaster with greater fiscal space, yet Maule's GDP trajectory shows no positive overshoot. New Zealand, despite higher sovereign indebtedness, channeled reconstruction capital primarily through its pre-existing insurance system rather than fiscal expansion. The magnitudes involved are telling: NZ's total insured cost reached NZ\$31~billion, of which more than NZ\$5~billion flowed from international reinsurers (Table~\ref{tab:quantcontrasts}), an exogenous capital injection that required no fiscal outlay. Chile's total insured cost was only US\$4--7~billion, with residential penetration below 20\%, leaving a US\$8.4~billion fiscal commitment as the primary reconstruction channel. Accordingly, Chile's change in public spending in the quake year (+10.44\%) exceeded New Zealand's (+6.24\%), directly reflecting this heavier reliance on the fiscal channel. The chronology in Table~\ref{tab:chronology} reinforces the point: NZ's EQC activated its claims process within days and appointed Fletcher Construction within weeks, whereas Chile's Plan de Reconstrucci{\'o}n was published six months after the earthquake. This contrast suggests that what matters for recovery dynamics is not necessarily raw fiscal capacity \textit{per se}, but the institutional and financial architecture through which reconstruction capital is mobilized. This observation does not imply that fiscal capacity is irrelevant, but rather that the \textit{speed and institutional channel} of financing may matter more for regional output than aggregate fiscal capacity.

\textbf{Legal system, property rights, and regulatory quality:} New Zealand scores higher on Legal System \& Property Rights (8.69 vs.\ 6.80) and Regulation (8.92 vs.\ 7.55). These national-level scores are consistent with the disaster-specific institutional evidence documented in Tables~\ref{tab:chronology} and~\ref{tab:regmilestones}: New Zealand enacted the Canterbury Earthquake Response and Recovery Act within 10~days of the September 2010 event, established CERA with statutory override powers within 35~days of the February 2011 shock, gazetted technical land categories and building consent fast-tracking within 12~months, and had the SCIRT public--private alliance and Land Use Recovery Plan operational within 24~months. Over CERA's lifespan, more than 60 Orders in Council were made to amend existing legislation for recovery purposes (Table~\ref{tab:regmilestones}). Chile, by contrast, made no comparable regulatory adaptations: reconstruction was managed through existing MINVU and MOP structures without a dedicated recovery entity, special legislative powers, or consent fast-tracking. The higher scores on Legal System \& Property Rights also align with the insurance channel: a well-functioning property rights regime is a \textit{precondition} for the complex contractual architecture of catastrophe insurance, reinsurance treaties, and the claims adjudication process through which more than 470,000 EQC and 168,000 private insurer claims were processed.

\textbf{Size of Government:} Chile's higher score (7.91 vs.\ 4.91) indicates a smaller government footprint relative to GDP (lower taxation, less government consumption, and fewer transfers). This is consistent with Chile's reliance on centralized but \textit{ad hoc} fiscal programs for reconstruction: the US\$8.4~billion commitment was channeled through existing ministry structures with no new institutional entities, whereas New Zealand's larger public sector provided the institutional infrastructure through which recovery was coordinated. Specifically, EQC (a Crown entity established in 1945 with NZ\$6.1~billion in accumulated reserves by 2010), CERA (a new government department with NZ\$4~billion in spending over five years), and the SCIRT alliance (a public--private arrangement coordinating three government entities and five construction firms on a NZ\$2.2~billion programme) were all creatures of New Zealand's comparatively larger state apparatus. The contrast illustrates that the relevant dimension is not government size in the abstract but rather the pre-existing institutional architecture that can be activated and adapted in a crisis. Chile's smaller government had correspondingly less pre-existing institutional capacity for disaster coordination (no public insurance scheme, no dedicated recovery authority) and the six-month lag to the Plan de Reconstrucci{\'o}n (Table~\ref{tab:chronology}) reflects the time required to build such capacity from scratch within existing ministerial structures.

To conclude, interpreting Table~\ref{tab:fiscal_inst} alongside the disaster-specific evidence in Tables~\ref{tab:chronology} to \ref{tab:regmilestones} reinforces the core interpretive point of Section~5: the divergent recovery trajectories are most consistently explained not by differential fiscal capacity (which actually favored Chile, at 11\% vs.\ 63\% debt-to-GDP) nor by broad macroeconomic stability, but by the pre-existing institutional architecture for disaster-specific financial mobilization, regulatory adaptation, and coordinated implementation. Chile's superior fiscal position did not translate into a comparable output response because the institutional channel through which spending was deployed (i.e., centralized ministerial programs without a dedicated recovery authority, or insurance-based external capital) was slower and less regionally targeted. The six-month lag to the Plan de Reconstrucci{\'o}n, the absence of regulatory override powers, and the limited community co-production all contrast with New Zealand's 35-day CERA establishment, over 60 Orders in Council adapting legislation, and the SCIRT public--private alliance.
